% Generated by scripts/build_reports.py; edit reviewed Markdown inputs.
% Compile with XeLaTeX.
% Options for packages loaded elsewhere
\PassOptionsToPackage{unicode}{hyperref}
\PassOptionsToPackage{hyphens}{url}
\PassOptionsToPackage{dvipsnames,svgnames,x11names}{xcolor}
\documentclass[
  11pt,
]{article}
\usepackage{xcolor}
\usepackage[margin=1in]{geometry}
\usepackage{amsmath,amssymb}
\setcounter{secnumdepth}{-\maxdimen} % remove section numbering
\usepackage{iftex}
\ifPDFTeX
  \usepackage[T1]{fontenc}
  \usepackage[utf8]{inputenc}
  \usepackage{textcomp} % provide euro and other symbols
\else % if luatex or xetex
  \usepackage{unicode-math} % this also loads fontspec
  \defaultfontfeatures{Scale=MatchLowercase}
  \defaultfontfeatures[\rmfamily]{Ligatures=TeX,Scale=1}
\fi
\usepackage{lmodern}
\ifPDFTeX\else
  % xetex/luatex font selection
\fi
% Use upquote if available, for straight quotes in verbatim environments
\IfFileExists{upquote.sty}{\usepackage{upquote}}{}
\IfFileExists{microtype.sty}{% use microtype if available
  \usepackage[]{microtype}
  \UseMicrotypeSet[protrusion]{basicmath} % disable protrusion for tt fonts
}{}
\makeatletter
\@ifundefined{KOMAClassName}{% if non-KOMA class
  \IfFileExists{parskip.sty}{%
    \usepackage{parskip}
  }{% else
    \setlength{\parindent}{0pt}
    \setlength{\parskip}{6pt plus 2pt minus 1pt}}
}{% if KOMA class
  \KOMAoptions{parskip=half}}
\makeatother
\usepackage{longtable,booktabs,array}
\usepackage{caption}
\captionsetup[table]{skip=6pt}
\newcounter{none} % for unnumbered tables
\usepackage{calc} % for calculating minipage widths
% Correct order of tables after \paragraph or \subparagraph
\usepackage{etoolbox}
\makeatletter
\patchcmd\longtable{\par}{\if@noskipsec\mbox{}\fi\par}{}{}
\makeatother
% Allow footnotes in longtable head/foot
\IfFileExists{footnotehyper.sty}{\usepackage{footnotehyper}}{\usepackage{footnote}}
\makesavenoteenv{longtable}
\setlength{\emergencystretch}{3em} % prevent overfull lines
\providecommand{\tightlist}{%
  \setlength{\itemsep}{0pt}\setlength{\parskip}{0pt}}
\usepackage{amsmath,amssymb}
\usepackage{xurl}
\usepackage{fancyhdr}
\pagestyle{fancy}
\fancyhf{}
\fancyhead[L]{\small Pure Thought | Revised research specifications}
\fancyfoot[L]{\small 2026-09-17}
\fancyfoot[R]{\small\thepage}
\setlength{\headheight}{15pt}
\setlength{\emergencystretch}{3em}
\allowdisplaybreaks
\usepackage{bookmark}
\IfFileExists{xurl.sty}{\usepackage{xurl}}{} % add URL line breaks if available
\urlstyle{same}
% fallback for those not using the hyperref driver hyperxmp:
\makeatletter
\@ifundefined{xmpquote}{\newcommand{\xmpquote}[1]{#1}}{}
\makeatother
\hypersetup{
  pdftitle={S08. Gapped boundaries from exact anyon data},
  pdfauthor={Pure Thought research portfolio},
  colorlinks=true,
  linkcolor={blue},
  filecolor={Maroon},
  citecolor={Blue},
  urlcolor={blue},
  pdfcreator={LaTeX via pandoc}}

\title{S08. Gapped boundaries from exact anyon data}
\author{Pure Thought research portfolio}
\date{Revised 17 September 2026}

\begin{document}
\maketitle
\thispagestyle{empty}
\clearpage

\section{S08. Gapped boundaries from exact anyon data}

\textbf{Source:} synthesis challenge 24;
\href{https://github.com/rick-stevens-ai/Pure-Thought/blob/main/NOTES/30-Most-Compelling-Pure-Thought-AI-Challenges.md}{original
synthesis}. This topic is distinct from, or adds a distinct objective
to, the canonical PRD catalog.
\href{https://github.com/rick-stevens-ai/Pure-Thought/blob/main/reviews/SYNTHESIS-CROSSWALK.md}{Crosswalk}

\subsection{Plain-language guide}

\subsubsection{The problem in everyday terms}

In some two-dimensional quantum phases, collective excitations behave
neither like ordinary bosons nor fermions. Their possible combinations
and exchanges restrict how the phase can end at an edge. This problem
asks which proposed ways of absorbing these excitations at a boundary
are mathematically consistent and can describe an edge with a nonzero
energy gap.

\subsubsection{Key terms}

\begin{itemize}
\item
  \textbf{Topological order} --- A type of quantum organization
  characterized by nonlocal structure and, here, anyon excitation data.
\item
  \textbf{Anyon} --- A two-dimensional excitation with exchange or
  braiding behavior beyond ordinary bosonic and fermionic possibilities.
\item
  \textbf{Fusion} --- The possible outcomes when excitations are
  combined.
\item
  \textbf{Braiding} --- The transformation associated with moving
  excitations around one another.
\item
  \textbf{Topological twist} --- The phase associated with a full
  rotation of an anyon.
\item
  \textbf{F and R data} --- Numbers or matrices specifying how fusion
  regrouping and exchanges act in chosen conventions.
\item
  \textbf{Condensate} --- A collection of excitations treated as
  absorbable into the boundary vacuum, together with consistent
  combination rules.
\item
  \textbf{Lagrangian algebra} --- The mathematical structure expressing
  a sufficiently complete compatible condensate for the boundary
  considered here.
\item
  \textbf{Gapped boundary} --- An edge whose excitations retain a
  positive energy threshold in the appropriate large-system limit.
\end{itemize}

\subsubsection{Why this matters}

Boundaries determine how a topological phase interfaces with its
surroundings and which excitations can disappear or remain
distinguishable there. Understanding these options supports the theory
of topological phases and the organization of quantum information in
such systems.

\subsubsection{What progress would mean}

Progress would verify a boundary candidate using complete consistency
data, or reject it for a precise algebraic reason. Matching a list of
fusion outcomes alone is insufficient. Turning the algebraic result into
a particular microscopic boundary Hamiltonian is an additional
construction.

\subsection{Revised core question}

For one fully specified two-dimensional topological order, can candidate
condensates and their gapped boundaries be verified from exact
categorical data?

\subsection{Precise scope, target and outputs}

Start with the toric-code modular category, explicit simple objects,
fusion, twists, F and R symbols and conventions. Enumerate candidate
condensable algebra objects within a declared multiplicity bound. Check
multiplication/unit, associativity, commutativity and separability
conditions, and the Lagrangian dimension condition for a boundary to
vacuum. For the abelian toric-code benchmark recover electric and
magnetic condensates. If constructing a lattice boundary, supply its
Hamiltonian and check commuting-projector and gap claims separately.

\subsection{Detailed critique of the synthesis framing}

Fusion rules and modular S,T tables alone do not in general specify or
prove all categorical coherence data. The synthesis also mixes 2D anyons
with 3D topological orders, which require different structures. A listed
condensate is not automatically a microscopic boundary realization. The
rewrite retains an exact categorical task, bounds the enumeration and
separates algebraic classification from lattice realization and
anomaly/inflow interpretation.

\subsection{Validation and milestones}

Baseline: verify the toric-code coherence data and two standard vacuum
boundaries. Strong: a second fully specified category or a completeness
proof within a bounded class. Invalid F/R symbols and a nonbosonic
condensate are rejection controls.

Use the shared
\href{https://github.com/rick-stevens-ai/Pure-Thought/blob/main/reviews/EVIDENCE-STANDARD.md}{evidence
standard}. Unsuccessful search is unresolved unless a complete exclusion
proof is supplied. These are proposed baseline specifications, not
completed scientific results.

\href{https://arxiv.org/abs/2504.19512}{Li, Yang and Dong, Gapped
Boundaries of Kitaev's Quantum Double Models} provides a primary example
of linking Lagrangian algebras to explicit lattice boundary terms.

\subsection{Displayed formulation and claim boundary}

\href{https://github.com/rick-stevens-ai/Pure-Thought/blob/main/supplemental/S08.md}{Supplemental
specification}.

\[A=\bigoplus_a n_a a,\qquad n_a\in\mathbb Z_{\ge0},\qquad
(\dim A)^2=\operatorname{Dim}\mathcal C.\]

For a vacuum boundary this Lagrangian dimension condition supplements,
rather than replaces, the unit, multiplication, associativity,
commutativity and separability checks using exact categorical coherence
data. In the toric-code category the reference condensates are
\(1\oplus e\) and \(1\oplus m\). A microscopic boundary Hamiltonian
needs its own gap proof.

\section{Evidence and execution standard}

This standard applies to the revised specifications. It is a proposed
contract for future research runs; this editorial revision does not
supply the scientific certificates themselves.

\subsection{What ``pure thought'' means here}

Use explicitly stated mathematical models and deductions without fitting
conclusions to experimental outcomes. Synthetic examples, exact model
inputs, analytic references, numerical search, and established software
are allowed. Physical constants, electronic-structure approximations,
constitutive laws, and empirical parameter sets must be identified when
used. A model-derived result is conditional on that model.

The original ban on legacy software conflicts with its proposed reliance
on NumPy, SciPy, CVXPY, nauty, and other libraries. Reimplementing every
optimizer and algebra system is not necessary for scientific
independence. Prefer a small independently checkable certificate and a
documented trusted base. A fresh implementation can be a project
objective, but must not be confused with a mathematical proof or novelty
claim.

\subsection{Four evidence levels}

{\def\LTcaptype{none} % do not increment counter
\begin{longtable}[]{@{}
  >{\raggedright\arraybackslash}p{(\linewidth - 4\tabcolsep) * \real{0.3333}}
  >{\raggedright\arraybackslash}p{(\linewidth - 4\tabcolsep) * \real{0.3333}}
  >{\raggedright\arraybackslash}p{(\linewidth - 4\tabcolsep) * \real{0.3333}}@{}}
\toprule\noalign{}
\begin{minipage}[b]{\linewidth}\raggedright
Label
\end{minipage} & \begin{minipage}[b]{\linewidth}\raggedright
What it establishes
\end{minipage} & \begin{minipage}[b]{\linewidth}\raggedright
What it does not establish
\end{minipage} \\
\midrule\noalign{}
\endhead
\bottomrule\noalign{}
\endlastfoot
Numerical evidence & A reproducible computation at stated precision,
sample size, and discretization & Exact truth, continuum coverage, or a
universal theorem \\
Finite-model certificate & An exact/validated statement about the
specified finite model or finite relaxation & An omitted tail,
thermodynamic limit, or physical realization \\
Conditional mathematical theorem & A proof over the stated domain with
every hypothesis and error controlled & Truth outside those assumptions
or a demonstrated physical system \\
Formal proof & A checked theorem in a named proof assistant, with
imported axioms/dependencies listed & Correct physical modeling merely
by virtue of formalization \\
\end{longtable}
}

Formalization is an optional extra deliverable unless the specific
project makes it mandatory. High precision, JSON output, solver status
and small residuals do not by themselves confer any of these proof
levels.

\subsection{Claim directions}

For minimization over a larger feasible set, the optimum is a lower
bound on the original minimum. A physical feasible point supplies an
upper bound. For maximization the directions reverse. Track feasibility
of the primal/dual witnesses separately from their objective values.

Infeasibility of an inner ansatz excludes only that ansatz.
Infeasibility of an outer relaxation can exclude the original problem if
the map into the relaxation is proved and all approximations preserve
necessity. Passing necessary tests means ``not excluded''; existence
needs a constructive sufficient argument.

A smallest observed example is not a minimality theorem. A complete
catalog must define the universe, equivalence relation, boundaries, and
proof that pruning misses nothing. A finite-size numerical trend does
not prove an asymptotic limit.

\subsection{Minimum run manifest}

Before calculation, record the problem ID and revision hash, exact
inputs and conventions, quantified domain, equivalence relation, target
observable, accepted evidence level, tolerances, and resource limits.
Record dependencies and precision/rounding mode. Where randomization is
used, preserve seeds, trial counts and raw sufficient statistics.
Resource limits are planning decisions, not evidence that a mathematical
solution cannot exist.

An output bundle should include:

\begin{itemize}
\tightlist
\item
  The statement being claimed, including all assumptions and
  quantifiers.
\item
  Construction or proof object, plus an independently runnable verifier
  and its trust assumptions.
\item
  Numerical/error ledger separating roundoff, discretization,
  truncation, optimization, sampling and modeling error where relevant.
\item
  Positive reference cases and negative controls that the verifier must
  reject.
\item
  Coverage: included domain, omitted domain, unresolved regions and
  termination reason.
\end{itemize}

These are requirements for future implementations, not a promise of
polynomial-size certificates for every problem.

\subsection{Checking specific certificate types}

\textbf{LP/SDP/SOS:} verify equality constraints, objective direction,
and cone feasibility using exact arithmetic or outward-rounded
intervals. Rationalizing a numerical solution can break PSD constraints.
Correct residuals with justified norm bounds; do not silently project
and reuse the old objective. Distinguish a proof of a polynomial
approximation from a proof for the original function.

\textbf{Interval numerics:} state function enclosures, domains and
outward rounding. A small root box needs existence and, if claimed,
uniqueness. Cover the complement for a global enumeration. Include
singularities and tails rather than discarding them numerically.

\textbf{SAT/SMT:} a solver's UNSAT label is not a certificate. Preserve
a proof supported by a checker and the proof that the encoding, symmetry
reduction and blocked solutions match the mathematical universe. DRAT
applies to an appropriate Boolean CNF proof, not automatically to every
real-arithmetic or SDP claim.

\textbf{Statistics:} specify the estimand, dependence assumptions,
confidence procedure and bias. Zero observed failures does not prove
zero failure probability. MCMC error bars based on estimated
autocorrelation are not generally rigorous mixing guarantees.

\subsection{Milestones and novelty}

Baseline means a working, independently validated result on a concrete
reference problem. Strong means a broader or sharper result with its
claimed evidence level. Research extension means a possible novel result
whose feasibility and significance need review. A negative theorem can
be valuable; an unsuccessful search can also be useful if its scope is
documented, but the two are not interchangeable.

Do not promise ``publication-quality'' in a fixed number of months.
Before claiming novelty, conduct a focused current literature comparison
against the exact theorem/model/metric. The sources in this revision
check key definitions and limitations; they are not an exhaustive
novelty audit across all 30 fields.

\end{document}
